\subsection{Case-Study Role and Model Contract}
\label{sec:methodology}

This section applies the review framework to one controlled stenosis example.
It identifies the model tier, native variables, closure choices, boundary
approximation, principal numerical method, observation operator, and evidence
contract before any result is interpreted. Dense operator formulas, secondary
solver surfaces, and reproducibility commands remain in the appendices unless
they are needed to understand a main-body claim.

\begin{center}
\fbox{\begin{minipage}{0.90\linewidth}
\small
\textbf{Solver-coordinate convention.} In the comparison record, the stored
solver state \texttt{A} is $a=R^2$, not physical circular area. The physical
map is Definition~\ref{def:physical-to-solver-map}:
$A_{\mathrm{phys}}=\pi a$, $Q_{\mathrm{phys}}=\pi q$, and
$\bar u=q/a=Q_{\mathrm{phys}}/A_{\mathrm{phys}}$.
\end{minipage}}
\end{center}

\begin{table}[!htb]
    \centering
    \scriptsize
    \caption{Case-study contract for the retained 23\% and 40\%
    stenosis-severity numerical evidence.
    The table states the minimum computational and interpretive commitments
    needed before the result tables are read. It is a reading contract, not an
    additional numerical result.}
    \label{tab:t1-case-study-contract}
    \begin{tabularx}{\textwidth}{@{}
        >{\raggedright\arraybackslash}p{0.20\textwidth}
        >{\raggedright\arraybackslash}p{0.39\textwidth}
        >{\raggedright\arraybackslash}X@{}}
        \toprule
        \textbf{Contract item}
            & \textbf{Main-body statement}
            & \textbf{Interpretation boundary} \\
        \midrule
        Case-study purpose
            & A controlled idealized stenosis example applies the review
            framework to a selected 1D area--flow model and two resolved
            velocity fields.
            & The chapter illustrates verification, reproducibility, and
            diagnostic cross-model velocity comparison; it is not a
            pressure-ratio study or patient-specific prediction. \\
        Resolved datasets and final-time match
            & The retained resolved datasets are the 23\% stenosis case from
            local dataset 77 and the 40\% stenosis case from local dataset 60.
            Both are compared at $T=0.9995$ s, with recorded 1D--3D final-time
            offsets of $5.47\times10^{-14}$ s.
            & The final timestamp is matched. Earlier transient-history
            matching is not carried by the retained record and would be needed
            for a stronger predictive comparison. \\
        Principal 1D computation
            & The main comparison uses the $R_{\max}$-normalized
            Canic-derived extended 1D member with Newtonian
            $\nu=0.04$ cm$^2$/s, parabolic profile closure, $N=400$ MUSCL
            finite-volume cells, Rusanov flux, and native SSPRK3 time
            integration.
            & Other finite-volume, DG, timestepper, and descriptor records are
            implementation-health or sensitivity context rather than the
            principal two-case comparison method. \\
        Observation operator
            & The 3D velocity field is reduced by the declared
            \texttt{CrossSectionQuadratureOperator}; the 1D state is observed as
            physical flow $Q_{1D}=\pi q$ and mean velocity
            $\bar u_{1D}=q/a$ on 200 axial sections. A second retained
            comparison applies the resolved displacement field before forming
            the same plane cuts.
            & The comparison is meaningful only for these common section
            observables. Radial-profile outputs are retained only as an audit
            record because their area-closure gate fails. \\
        Numerical admissibility gates
            & The main text reports manufactured-solution behavior,
            output-operator verification, geometry-rest drift, subcritical
            boundary diagnostics, CFL, positivity, and finite-volume
            bookkeeping.
            & These records support run admissibility and numerical
            interpretation. They do not by themselves identify a physical
            prediction target. \\
        Matching limits
            & The analytic reference profile is shared and the static cut-area
            check is reported for both reference and deformed coordinates. A
            quasi-static membrane-FSI comparator is also retained for the
            radial wall displacement solve.
            & Velocity differences cannot be attributed uniquely to a specific
            closure, wall, boundary, or geometry mechanism because the retained
            resolved record still lacks matched boundary histories and full
            transient wall history. \\
        Retained outputs
            & The retained numerical evidence is the MMS/rest-state record,
            cut-area check, axial flow behavior, flow/area bookkeeping,
            section-mean velocity discrepancy summary, coordinate-mode
            comparison, quasi-static membrane-FSI wall summary,
            production-grid sensitivity, and largest section discrepancies.
            & These are descriptive final-time discrepancy and sensitivity
            records for the declared two-case comparison. \\
        Excluded radial/profile outputs
            & Secondary radial-profile rows were regenerated for selected
            axial slices and audited.
            & The raw rows reproduce the generated radial summary, but the
            area-closure gate fails because the radial slices do not coincide
            with the retained section-target set. \\
        \bottomrule
    \end{tabularx}
\end{table}

\subsubsection{Model Contract and Native Variables}
\label{subsec:methodology-problem-definition}
\label{subsec:selected-1d-forward-model}
\label{subsec:cross-sectional-variables-governing-equations}
\label{subsubsec:cross-sectional-reduction-selected-1d-model}

The computational problem is a single idealized stenotic vessel represented by a
selected one-dimensional area--flow model. The model assumes a distinguished
axial direction and cross-sectionally averaged pressure and velocity, retaining
distributed pulse-wave and pressure--flow structure while discarding secondary
flow, separation, recirculation, and resolved wall traction. The reduction
follows standard compliant-artery and dimension-reduced modeling arguments
\cite[Sec.~2, pp.~253--257]{FormaggiaEtAl2003OneDimensionalBloodFlow}
\cite[Ch.~3, Secs.~3.2 and~3.6]{KopplHelmig2023DimensionReducedModeling}
\cite[Sec.~3, pp.~24--26]{ShiEtAl2011BloodFlowModelReview}.

\begin{assumption}[Straight single-vessel geometry]
    \label{ass:1d-straight-vessel-geometry}
    The selected 1D model assumes a vessel of length $L$ aligned with the
    $z$-axis and a circular cross-section with radius $R(z,t)$. For
    $0<z<L$, $S(z,t)$ is the interior cross-section used for averaging, not
    a boundary component of the open lumen.
\end{assumption}

\begin{definition}[Section-averaged fields and observed pressure]
    \label{def:section-averaged-1d-fields}
    The retained reduced kinematic variables and observed pressure
    representative are
    \begin{flalign*}
        \quad A_{\mathrm{phys}}(z,t) &\defined \int_{S(z,t)} 1 \, dS, &&
            \text{physical cross-sectional area}, \\
        \quad Q_{\mathrm{phys}}(z,t) &\defined \int_{S(z,t)} u_z \, dS, &&
            \text{physical volumetric flow rate}, \\
        \quad \bar u(z,t) &\defined
            \frac{Q_{\mathrm{phys}}(z,t)}{A_{\mathrm{phys}}(z,t)}, &&
            \text{mean axial velocity}, \\
        \quad \bar p_{3D,\mathrm g}(z,t) &\defined
            \frac{1}{A_{\mathrm{phys}}(z,t)}
            \int_{S(z,t)} p_{3D,\mathrm g} \, dS, &&
            \text{observed 3D cross-sectional mean gauge pressure}.
    \end{flalign*}
    Here $p_{3D,\mathrm g}$ is continuum pressure in the wall-law gauge frame.
    The observed representative $\bar p_{3D,\mathrm g}$ is distinct from the
    reduced wall-law pressure variable introduced below.
\end{definition}

\begin{definition}[Physical-to-solver map for the implemented 1D state]
    \label{def:physical-to-solver-map}
    The continuum 1D model uses physical area $A_{\mathrm{phys}}$ and
    physical flow $Q_{\mathrm{phys}}$. The implemented radius-squared solver
    instead stores
    \begin{flalign*}
        \quad
        a=R^2,\qquad q=\frac{Q_{\mathrm{phys}}}{\pi}.
        &&
    \end{flalign*}
    Thus
    \begin{flalign*}
        \quad
        A_{\mathrm{phys}}=\pi a,\qquad
        Q_{\mathrm{phys}}=\pi q,\qquad
        \bar u=\frac{Q_{\mathrm{phys}}}{A_{\mathrm{phys}}}=\frac{q}{a}.
        &&
    \end{flalign*}
    With $K=Eh/(1-\sigma^2)$ and the Canic conservative-form convention
    $R_0^*=R_{\max}$, the selected evolution wall law gives the solver elastic
    flux contribution recorded in Appendix~\ref{app:num-current-solver-operator}:
    \begin{flalign*}
        \quad
        \Psi_h(a)=\frac{K}{3\rho R_{\max}^2}a^{3/2},
        \qquad
        \pd_a\Psi_h(a)=\frac{K}{2\rho R_{\max}^2}\sqrt a.
        &&
    \end{flalign*}
    This map is the convention used by the
    Section~\ref{sec:resolved-3d-comparison} velocity and flow comparisons.
\end{definition}

\subsubsection{Closure Choices: Geometry and Wall Law}
\label{subsec:stenosis-geometry-implemented-coordinates}
\label{subsec:selected-wall-law-closure-model}

\begin{definition}[Stenosis reference geometry]
    \label{def:1d-stenosis-reference-geometry}
    A stenosis is encoded through the reference radius or reference area. Let
    $R_{\mathrm{base}}(z)$ be the nominal healthy radius and $s(z)\in[0,1)$
    a fractional radius-reduction profile. Then
    \begin{flalign*}
        \quad
        R_0(z) &= R_{\mathrm{base}}(z)\big(1-s(z)\big),
        &
        A_0(z) &= \pi R_0(z)^2.
        \numberthis\label{eq:1d-stenosis-area} &&
    \end{flalign*}
    The baseline idealized stenotic-vessel profile is the smooth asymmetric
    kernel
    \begin{flalign*}
        \quad
        g(z) &= z-z_s+\kappa
        \exp\left(-\frac{(z-z_a)^2}{2\sigma_a^2}\right),
        &
        \psi(z) &= \exp\left(-\lambda g(z)^4\right),
        \numberthis\label{eq:1d-stenosis-kernel} && \\
        \quad
        s(z) &= s_{\max}\psi(z), \qquad 0\leq s_{\max}<1.
        \numberthis\label{eq:1d-stenosis-profile} &&
    \end{flalign*}
    In the default centimeter-scale simulation record, $z_a=2.5\,\mathrm{cm}$,
    $z_s=3.4\,\mathrm{cm}$, $\sigma_a=1\,\mathrm{cm}$,
    $\kappa=0.95\,\mathrm{cm}$, and
    $\lambda=50\,\mathrm{cm}^{-4}$.
    This geometry is the baseline idealized simulation geometry for the worked
    example.
\end{definition}

\begin{remark}[Stenosis-profile and variable-radius boundary]
    \label{rmk:1d-stenosis-profile-scope}
    The profile in Eq.~\ref{eq:1d-stenosis-profile} is $C^\infty$ and
    rapidly localized rather than compactly supported. This smooth idealized
    geometry is the baseline for the reduced simulations in this report and is
    also useful for fixed-wall 3D studies: the same analytic radius profile
    supplies controlled derivatives, wall normals, severity variation, and
    repeatable mesh-refinement families without introducing segmentation noise.
    Canic et al.\ show that standard 1D variable-geometry models may miss
    stenotic pressure and velocity trends unless additional variable-radius
    terms are included \cite[Abstract; Secs.~1--2]{CanicEtAl2024Extended1DStenosis}.
    Those corrections are included by the \texttt{canic-extended-1d}
    forward-model descriptor; the \texttt{classical-1d-no-slip} descriptor is
    the parabolic-profile baseline that retains the same selected wall law but
    omits the Canic variable-radius corrections.
\end{remark}

\begin{definition}[Classical parabolic-profile 1D baseline]
    \label{def:classical-no-slip-1d-baseline}
    The classical baseline used as a model-family member is a
    parabolic-profile reduced model. A no-slip wall trace is compatible with a
    parabolic profile for fully developed Newtonian flow in a straight circular
    tube, but no-slip alone does not imply that profile in an unsteady stenotic
    vessel. The parabolic profile is therefore an independent cross-sectional
    closure with $\alpha_{\mathcal P}=4/3$ and $g_{\mathcal P}=4$. In the
    implementation manifests this member is recorded under the descriptor
    \texttt{classical-1d-no-slip}; in the mathematical discussion it is
    the classical parabolic-profile baseline. It keeps the same reduced state
    variables and selected Canic-derived pressure--area wall law as the
    extended 1D framework, but omits the Canic effective-alpha variable-radius
    correction. The descriptor \texttt{canic-extended-1d} names the
    $R_{\max}$-normalized Canic-derived extended-model member.
\end{definition}

\begin{definition}[Selected Canic--Koiter pressure--area wall law]
    \label{def:1d-elastic-wall-law}
    The selected elastic wall closure is the Canic--Koiter thin-wall
    pressure--area law for the reduced wall-law gauge pressure
    $p_{\mathcal W,\mathrm g}$. In conventional area notation it may be written
    \begin{flalign*}
        \quad p_{\mathcal W,\mathrm g}(A,z,t)
        =
        p_{\mathrm{ext}}(z,t)
        + \frac{\beta(z)}{A_0(z)}
            \left(\sqrt{A(z,t)}-\sqrt{A_0(z)}\right).
        \numberthis\label{eq:1d-wall-law} &&
    \end{flalign*}
    Here $A_0(z)$ is the reference area, $p_{\mathrm{ext}}$ is the external
    pressure, and $\beta(z)$ is a wall-stiffness parameter. In the
    radius-squared coordinate $a=R^2$ used by the implemented finite-volume
    evolution, the selected continuous source-balanced member uses the Canic
    conservative-form reference radius $R_0^*=R_{\max}$:
    \begin{flalign*}
        \quad
        p_{1,\mathrm g}(a,z,t)
        =
        p_{\mathrm{ext}}(z,t)
        +
        \frac{K}{R_{\max}^2}
        \left(\sqrt a-R_0(z)\right),
        \qquad
        K=\frac{Eh}{1-\sigma^2}. &&
    \end{flalign*}
    With $A=\pi a$ and $A_0=\pi R_0^2$, the conventional law in
    Eq.~\ref{eq:1d-wall-law} reduces to this selected member only when
    \begin{flalign*}
        \quad
        \beta(z)=\sqrt{\pi}\,K\,\frac{R_0(z)^2}{R_{\max}^2}. &&
    \end{flalign*}
    Thus the implemented evolution uses a constant reference-radius
    denominator, not the local $R_0^{-2}$ denominator. In this report
    $R_{\max}$ is the healthy reference-radius parameter used by the solver;
    for the default stenosis profile it equals $\max_zR_0(z)$. The selected
    comparison runs use the cgs numerical-record values
    $\rho=1.055\,\mathrm{g}/\mathrm{cm}^3$,
    $E=5.02\times10^6\,\mathrm{dyn}/\mathrm{cm}^2$,
    $h=0.06\,\mathrm{cm}$, $\sigma=0.5$,
    $K=4.016\times10^5\,\mathrm{dyn}/\mathrm{cm}$,
    $R_{\max}=0.18\,\mathrm{cm}$, and
    $p_{\mathrm{ext}}\equiv0$.
    This is the wall law paired with the implemented elastic potential and
    wall/geometry source below. Legacy pressure-output helpers in the
    implementation still report a local-denominator pressure diagnostic through
    \path{pressure(result, params)}; benchmark pressure extrema and
    pressure-diagnostic bars refer to that diagnostic, not to the selected
    evolution wall law. The 1D model identifies $p_{\mathcal W,\mathrm g}$ with
    a cross-sectional pressure representative as a closure assumption and does
    not resolve radial pressure variation.
\end{definition}

\subsubsection{Reduced Balance Law and Rest-State Contract}
\label{subsec:conservative-form-source-terms}

The selected 1D model is an area--flow balance law closed by the profile,
rheology, wall, geometry, and boundary choices stated above. In review terms,
the important point is that the model retains only axial area and flow while
moving transverse velocity, wall response, friction, rheology, and
variable-radius effects into closures. Appendix~\ref{app:num-current-solver-operator}
gives the corresponding semi-discrete operator in solver coordinates, including
the $R_{\max}$-normalized wall potential, wall/geometry source, friction term,
effective-alpha correction, and locally frozen $p_2$ derivative. The main-body
case study uses that appendix as the implementation authority and records here
only the model contract needed to interpret the evidence.

The continuous selected member has a rest state that matters for the evidence
standard. Under constant external pressure and compatible zero-flow boundary
data, the reference geometry $a(z)=R_0(z)^2$ with $q(z)=0$ should remain at
rest in the continuous source-balanced model. This property is the target of
the geometry-rest verification test in Section~\ref{subsec:rest-state-drift}.

\begin{proposition}[Rest-state compatibility of the selected wall source]
    \label{prop:selected-rest-state-compatibility}
    Suppose $p_{\mathrm{ext}}$ is constant in the gauge convention and the
    boundary data are compatible with $q=0$. Then
    $a(z)=R_0(z)^2$ and $q(z)=0$ are an exact steady state of the selected
    continuous source-balanced evolution law.
\end{proposition}

\begin{proof}
    Using the operator definitions in Appendix~\ref{app:num-current-solver-operator},
    the mass equation gives $\pd_z q=0$. In the momentum equation, all terms
    proportional to $q$ vanish, so
    $S_{\mathrm{fric}}=S_{\alpha}=S_{p_2}=0$. At $a=R_0^2$,
    \begin{flalign*}
        \quad
        \pd_z\Psi_h(R_0^2)
        &=
        \pd_z\left(\frac{K}{3\rho R_{\max}^2}R_0^3\right)
        =
        \frac{K}{\rho R_{\max}^2}R_0^2R_0'
        =
        S_{\mathrm{wall}}(R_0^2,z). &&
    \end{flalign*}
    Thus the elastic flux derivative is exactly canceled by the wall/geometry
    source, and $\pd_t q=0$.
\end{proof}

Continuous compatibility does not imply discrete preservation: the current
MUSCL/Rusanov method exhibits the reported non-well-balanced geometry-rest
drift.

\subsubsection{Source-to-Implementation Differences}
\label{subsec:source-to-implementation-variant}

The selected implementation is a documented Canic-derived extended model, not
a literal reproduction of every source-paper parameter and normalization in
\cite{CanicEtAl2024Extended1DStenosis}. The report therefore uses the phrase
\emph{$R_{\max}$-normalized Canic-derived extended model} when source fidelity
matters. Table~\ref{tab:source-to-implementation-model} records the relevant
source-to-implementation choices. The diagnostic extended pressure correction is
\begin{flalign*}
    \quad
    p_{2,\mathrm g}(a,q,z,t)
    =
    g_{\mathcal P}\rho\nu_{\mathrm{eff}}^{\mathcal R}
    \frac{q}{a}\frac{R_0'(z)}{R_0(z)}. &&
\end{flalign*}
This $p_2$ term is a variable-radius viscous correction, not a replacement wall
law. In generalized-Newtonian descriptor runs, the implemented analytic $p_2$
derivative uses the local value of $\nu_{\mathrm{eff}}^{\mathcal R}$ and does
not differentiate that effective-viscosity factor; it is therefore a locally
frozen-viscosity computational closure for this term.

\begin{table}[!htb]
    \centering
    \scriptsize
    \caption{Source-to-implementation map for the selected 1D model.}
    \label{tab:source-to-implementation-model}
    \begin{tabularx}{\textwidth}{@{}
        >{\raggedright\arraybackslash}p{0.24\textwidth}
        >{\raggedright\arraybackslash}p{0.34\textwidth}
        >{\raggedright\arraybackslash}X@{}}
        \toprule
        \textbf{Feature} & \textbf{Canic et al. source form} & \textbf{Implemented report form} \\
        \midrule
        Pressure denominator
            & Elastic pressure written with a local $R_0(z)^{-2}$ denominator.
            & Evolution wall potential and source use the constant
            $R_{\max}^{-2}$ normalization specified in
            Appendix~\ref{app:num-current-solver-operator}. \\
        Profile parameters
            & Source simulations state $\gamma=9$ and $\alpha=1.1$.
            & Principal runs use the parabolic profile
            $\alpha_{\mathcal P}=4/3$, $g_{\mathcal P}=4$; power-profile runs
            are sensitivity/diagnostic records. \\
        Friction coefficient
            & Profile-dependent wall/friction closure in the reduced model.
            & Implemented through the declared velocity-profile shear factor
            and rheology descriptor. \\
        $p_2$ differentiation
            & Variable-radius pressure correction differentiated in the model.
            & Locally frozen-viscosity treatment for the $p_2$ derivative. \\
        External pressure
            & External pressure term may be included in the formulation.
            & $p_{\mathrm{ext}}\equiv0$ in the reported runs. \\
        Coordinates
            & Physical area and flow notation in the source paper.
            & Solver state is $(a,q)$ with
            $A_{\mathrm{phys}}=\pi a$, $Q_{\mathrm{phys}}=\pi q$. \\
        Boundary rule
            & Boundary conditions are model/problem dependent.
            & Fixed-area characteristic approximation applied with persisted
            subcritical-regime diagnostics. \\
        Discrete well balancing
            & Continuous rest equilibrium is a model-level property.
            & Rusanov/MUSCL is not exactly well balanced for spatially varying
            $R_0$; Section~\ref{sec:numerical-verification} reports
            zero-forcing rest-state drift. \\
        \bottomrule
    \end{tabularx}
\end{table}

\subsubsection{Boundary Approximation and Subcritical Condition}
\label{subsec:initial-boundary-conditions}

The comparison runs start from the geometry-rest state and then impose the
finite-volume boundary approximation stated below.

\begin{definition}[Single-vessel boundary data]
    \label{def:1d-single-vessel-boundary-data}
    Generic single-vessel models may prescribe inlet flow or velocity and an
    outlet pressure or terminal relation; these alternatives are model-family
    context rather than the reported boundary rule. The realized finite-volume
    run prescribes an inlet solver flow derived from a nominal reference-area
    mean velocity,
    \begin{flalign*}
        \quad
        q_{\mathrm{in}}&=R_0(0)^2\bar u_{\mathrm{in}}, &&
    \end{flalign*}
    solves the inlet ghost area from the approximate outgoing characteristic
    invariant, fixes the outlet ghost coordinate
    $a_{\mathrm{out}}=R_0(L)^2$, and obtains the outlet ghost flow from the
    corresponding approximate invariant. Hence the imposed inlet condition is
    a prescribed reference-area flow, not a fixed realized mean velocity
    $q_{\mathrm{in}}/a_{\mathrm{in}}$, and the outlet condition fixes the
    elastic reference pressure only for the selected wall component.
\end{definition}

\begin{definition}[Characteristic speeds for the reduced elastic subsystem]
    \label{def:1d-characteristic-speeds}
    Let $\lambda_\pm$ denote the two characteristic speeds of the
    positive-area elastic 1D subsystem obtained from the homogeneous
    $(A,Q)$ balance law. Their formula and assumptions are given in
    Appendix~\ref{app:num-elastic-potential-form}.
\end{definition}

\begin{assumption}[Subcritical boundary regime]
    \label{ass:1d-subcritical-boundary-regime}
    The baseline boundary treatment is used only in the subcritical regime
    $\lambda_-<0<\lambda_+$, where one characteristic enters through the inlet
    and one characteristic enters through the outlet.
\end{assumption}

Under this regime, a finite-volume scheme constrains the incoming information at
each boundary and leaves the outgoing state determined by the interior solve.
The finite-volume boundary treatment in
Appendix~\ref{app:num-boundary-state-realization} uses an $\alpha=1$,
fixed-area characteristic approximation as the boundary invariant for the
variable-$\alpha_{\mathrm{eff}}$ solver. It is therefore an implemented boundary
approximation, not an exact invariant derivation for the full variable-radius
balance law.

\subsubsection{Principal MUSCL/Rusanov/SSPRK3 Method}
\label{subsec:muscl-ssprk3-discretization}

Only the principal comparison method belongs in the main narrative. The reported
two-case severity comparison uses the finite-volume method specified in
Appendix~\ref{app:num-current-solver-operator}: the radius-squared state
$(a,q)$ is advanced on a one-dimensional grid with MUSCL reconstruction,
minmod-limited conservative states, Rusanov flux, source splitting for the
wall/geometry and variable-radius terms, and native SSPRK3 time integration.
The principal numerical comparison uses $N=400$ cells on $0\le z\le6$ cm, a
timestep cap of $10^{-5}$ s, and a CFL cap of 0.45. The optional
output-sensitivity record reruns the same comparison protocol at
$N=200,400,800$ cells and is interpreted only as a production-grid sensitivity
check. The boundary-state approximation is the fixed-area characteristic
contract recorded in
Appendix~\ref{app:num-boundary-state-realization}. Other finite-volume, modal-DG,
native time-stepper, and SciML/OrdinaryDiffEq surfaces are implementation-check
or sensitivity context summarized in Table~\ref{tab:implemented-discretization-surface};
they are not the principal two-case comparison method.

\subsubsection{Plane--Tetrahedron Observation Operator}
\label{subsec:cross-sectional-velocity-reconstruction}
\label{subsec:plane-tetrahedron-comparison-operator}

The 1D state directly supplies a section mean,
$\bar u_{1D}(z,t)=q(z,t)/a(z,t)$, and physical flow
$Q_{1D}(z,t)=\pi q(z,t)$. The implementation can also emit radial-profile
rows using the same cut triangles. Those rows are audited in this manuscript,
but no radial-profile figure or localization claim is retained because the
radial area-closure gate fails.

The resolved-data loader also accepts pressure and displacement XDMF/HDF5
companions for a velocity field. When the comparison is explicitly run in
deformed-coordinate mode, the observation operator applies the node-centered
displacement before forming the plane cuts. Both the reference-coordinate and
deformed-coordinate section comparisons are retained in
Section~\ref{subsec:deformed-coordinate-comparison}.

For a target plane $S_j=\Omega_{\mathrm{3D}}\cap\{z=z_j\}$, the default
\texttt{CrossSectionQuadratureOperator} intersects every tetrahedron with
$S_j$, linearly interpolates axial velocity on the cut edges, sorts the cut
polygon vertices in the plane, and triangulates the polygon by a centroid fan.
Let $\mathcal Q_j$ denote the resulting cut triangles. The reported 3D area,
flow, and mean velocity are
\begin{flalign*}
    \quad
    A_{3D,j}
    &=
    \sum_{\Delta\in\mathcal Q_j}|\Delta|,
    &
    Q_{3D,j}
    &=
    \sum_{\Delta\in\mathcal Q_j}|\Delta|\,\bar u_{z,\Delta},
    &
    \bar u_{3D,j}
    &=
    Q_{3D,j}/A_{3D,j}. &&
\end{flalign*}
For a cut triangle $\Delta$ whose vertices carry reconstructed axial
velocities $u_{z,1}$, $u_{z,2}$, and $u_{z,3}$,
\begin{flalign*}
    \quad
    \bar u_{z,\Delta}
    =
    \frac{u_{z,1}+u_{z,2}+u_{z,3}}{3}. &&
\end{flalign*}
The corresponding 1D quantities use physical flow, not the raw solver flow
coordinate:
\begin{flalign*}
    \quad
    Q_{1D,j}
    =
    \pi q(z_j,t),
    \qquad
    \bar u_{1D,j}
    =
    q(z_j,t)/a(z_j,t). &&
\end{flalign*}
Quarantined radial rows are not part of the retained comparison metrics.

\subsubsection{Comparison Design and Discrepancy Metrics}
\label{subsec:numerical-experiments-metrics-provenance}

The resolved-velocity comparison evaluates the selected 1D solution against two
available resolved 3D velocity comparison datasets at the matched final-time
sample $T=0.9995$ s. The retained cases are described by their 23\% and 40\%
stenosis severities, with local dataset identifiers recorded only for
provenance. The comparison is a diagnostic cross-model velocity comparison
under the declared observation operator.
The compact reading contract is stated at the start of this section.
Regeneration commands, local source identifiers, raw-input policy, and static
asset ownership are recorded in
Appendix~\ref{app:code-and-ai-use}.

Velocity and flow differences are reported as sample discrepancy measures over
valid section targets. With
$d_{u,j}=\bar u_{1D,j}-\bar u_{3D,j}$ and
$d_{Q,j}=Q_{1D,j}-Q_{3D,j}$, the retained section metrics are the signed mean
bias, mean absolute discrepancy, RMS discrepancy, maximum absolute discrepancy,
and relative RMS discrepancy,
\begin{flalign*}
    \quad
    \overline d_u
    &=
    \frac{1}{n}\sum_j d_{u,j},
    &
    D_{u,1}
    &=
    \frac{1}{n}\sum_j |d_{u,j}|,
    &
    D_{u,2}
    &=
    \left(\frac{1}{n}\sum_j d_{u,j}^2\right)^{1/2}, &&\\
    \quad
    D_{u,\infty}
    &=
    \max_j |d_{u,j}|,
    &
    D_{u,2}^{\mathrm{rel}}
    &=
    \frac{\left(\sum_j d_{u,j}^2\right)^{1/2}}
         {\left(\sum_j\bar u_{3D,j}^2\right)^{1/2}}, &&\\
    \quad
    \overline d_Q
    &=
    \frac{1}{n}\sum_j d_{Q,j},
    &
    D_{Q,1}
    &=
    \frac{1}{n}\sum_j |d_{Q,j}|,
    &
    D_{Q,2}
    &=
    \left(\frac{1}{n}\sum_j d_{Q,j}^2\right)^{1/2}. &&
\end{flalign*}
